
\medskip
\phantomsection\label{def:sga3-VIB-7.2.2}%
\noindent\textbf{Definition 7.2.2.---}
Let \(G\) be a \(k\)-group.
\begin{enumerate}
\item[\textup{a)}]\phantomsection\label{def:sga3-VIB-7.2.2-a}
  Let \(A\) and \(B\) be two geometrically reduced
  \(k\)-group subschemes of \(G\). The \emph{commutator group
  subscheme of \(A\) and \(B\) in \(G\)}, denoted
  \[
    (A,B)_G
    \qquad\text{or simply}\qquad
    (A,B),
  \]
  is the closed group subscheme of \(G\) generated by the
  morphism
  \phantomsection\label{form:sga3-VIB-7.2.2-commutator}%
  \[
    \nu:A\times_kB\longrightarrow G
  \]
  defined, for every \(k\)-scheme \(S\), by
  \[
    (a,b)\longmapsto aba^{-1}b^{-1},
    \qquad
    a\in A(S),\quad b\in B(S).
  \]

\item[\textup{b)}]\phantomsection\label{def:sga3-VIB-7.2.2-b}
  Suppose that \(G\) is geometrically reduced over \(k\). The
  \emph{derived group of \(G\)}, denoted \(\mathfrak{D}(G)\),%
  \footnote{\label{note:sga3-VIB-73}Editor's note: the notation
  \(\operatorname{Der}(G)\) of the original has been changed to
  avoid any risk of confusion with a space of derivations.}
  is the group
  \[
    \mathfrak{D}(G)=(G,G).
  \]
\end{enumerate}

\noindent\emph{N.B.}
The group \(G\) is commutative if and only if \(\mathfrak{D}(G)\) is the
unit \(k\)-group.

\medskip
\phantomsection\label{recall:sga3-VIB-7.3.0}%
\noindent\textbf{Reminders \SGARefLink{recall:sga3-VIB-7.3.0}{7.3.0}.---}%
\footnote{\label{note:sga3-VIB-74}Editor's note: these reminders
have been added, and the statement of \SGARefLink{cor:sga3-VIB-7.3}{7.3} has consequently been
modified and expanded.}
Let \(\varphi:Y\to Z\) be a morphism of \(S\)-schemes. Recall
that the image presheaf \(\varphi(Y)\) is the \(S\)-functor
which assigns to every \(S'\) over \(S\) the subset
\[
  \varphi\bigl(Y(S')\bigr)\subset Z(S').
\]
Suppose that \(T\) is a subscheme of \(Z\). If the inclusion of
presheaves
\[
  \varphi(Y)\hookrightarrow Z
\]
factors through \(T\), that is, if
\[
  \varphi\circ h\in T(S')
\]
for every \(S\)-morphism \(h:S'\to Y\), then, set-theoretically,
\(\varphi(Y)\subset T\): take \(h=\operatorname{id}_Y\).
Conversely, if \(Y\) is reduced and
\(\varphi(Y)\subset T\) set-theoretically, then \(\varphi\)
factors through \(T\).

Recall also that, by \hyperref[prop:sga3-VIB-6.2]{6.2}, if
\(H\) is a closed \(k\)-group subscheme of \(G\), then
\[
  \underline{\operatorname{Centr}}_G H
  \qquad\text{and}\qquad
  \underline{\operatorname{Norm}}_G H
\]
are representable by closed \(k\)-group subschemes of \(G\).

\medskip
\phantomsection\label{cor:sga3-VIB-7.3}%
\noindent\textbf{Corollary 7.3.---}
\emph{Let \(G\) be a \(k\)-group, let \(X\) be a geometrically
reduced \(k\)-scheme, and let \(f:X\to G\) be a
\(k\)-morphism.}
\begin{enumerate}
\item[\textup{\SGAInlineRefTarget{cor:sga3-VIB-7.3:item:i}(i)}]\phantomsection\label{concl:sga3-VIB-7.3-i}
  \emph{Let \(S\) be a \(k\)-scheme and let \(u\) be an
  endomorphism of the \(S\)-group \(G_S\).}
  \begin{enumerate}
  \item[\textup{(a)}]\phantomsection\label{concl:sga3-VIB-7.3-i-a}
    \emph{If}
    \[
      u\bigl(f_S(X_S)\bigr)\subset\Gamma_G(f)_S
    \]
    \emph{set-theoretically, then the morphism}
    \[
      u:\Gamma_G(f)_S\longrightarrow G_S
    \]
    \emph{factors through \(\Gamma_G(f)_S\).}

  \item[\textup{\SGAInlineRefTarget{cor:sga3-VIB-7.3:item:b}(b)}]\phantomsection\label{concl:sga3-VIB-7.3-i-b}
    \emph{If \(u\) is an automorphism of the \(S\)-group \(G_S\)
    and}
    \[
      u\bigl(f_S(X_S)\bigr)=f_S(X_S)
    \]
    \emph{set-theoretically, then \(u\) induces an automorphism
    of \(\Gamma_G(f)_S\). In particular, if \(g\in G(S)\)
    satisfies}
    \[
      \operatorname{int}(g)\bigl(f_S(X_S)\bigr)=f_S(X_S)
    \]
    \emph{set-theoretically, then}
    \[
      g\in
      \underline{\operatorname{Norm}}_{G_S}
        \bigl(\Gamma_G(f)_S\bigr)(S).
    \]

  \item[\textup{\SGAInlineRefTarget{cor:sga3-VIB-7.3:item:c}(c)}]\phantomsection\label{concl:sga3-VIB-7.3-i-c}
    \emph{If \(u\circ f_S=f_S\), then the restriction of \(u\)
    to the subgroup \(\Gamma_G(f)_S\) of \(G_S\) is the identity.
    In particular, if \(g\in G(S)\) satisfies}
    \[
      \operatorname{int}(g)\circ f_S=f_S,
    \]
    \emph{then}
    \[
      g\in
      \underline{\operatorname{Centr}}_{G_S}
        \bigl(\Gamma_G(f)_S\bigr)(S).
    \]
  \end{enumerate}

\item[\textup{\SGAInlineRefTarget{cor:sga3-VIB-7.3:item:ii}(ii)}]\phantomsection\label{concl:sga3-VIB-7.3-ii}
  \emph{Let \(H\) be a group subscheme of \(G\). Then \(H\)
  centralizes (respectively, normalizes) \(\Gamma_G(f)\) if and
  only if, for every \(k\)-scheme \(S\) and every \(h\in H(S)\),}
  \[
    \operatorname{int}(h)\circ f_S=f_S
  \]
  \emph{(respectively,}
  \[
    \operatorname{int}(h)\bigl(f_S(X_S)\bigr)
      \subset\Gamma_G(f)_S
  \]
  \emph{set-theoretically). Equivalently, for every
  \(S'\to S\) and every \(x\in X(S')\), the elements \(h_{S'}\)
  and \(f(x)\) of \(G(S')\) commute (respectively,}
  \[
    h_{S'}f(x)h_{S'}^{-1}\in\Gamma_G(f)(S').
  \]
  \emph{).}

\item[\textup{\SGAInlineRefTarget{cor:sga3-VIB-7.3:item:iii}(iii)}]\phantomsection\label{concl:sga3-VIB-7.3-iii}
  \emph{In particular, let \(Y\) be a second geometrically
  reduced \(k\)-scheme and let \(\varphi:Y\to G\) be a
  \(k\)-morphism. Suppose that, for every \(k\)-scheme \(S'\),
  every \(x\in X(S')\), and every \(y\in Y(S')\), the elements
  \(\varphi(y)\) and \(f(x)\) of \(G(S')\) commute
  (respectively,}
  \[
    \operatorname{int}\bigl(\varphi(y)\bigr)\bigl(f(x)\bigr)
      \in\Gamma_G(f)(S').
  \]
  \emph{).}
  \emph{Then \(\Gamma_G(\varphi)\) is a \(k\)-subgroup of}
  \[
    \underline{\operatorname{Centr}}_G\bigl(\Gamma_G(f)\bigr)
    \qquad\text{\emph{(respectively, of)}}\qquad
    \underline{\operatorname{Norm}}_G\bigl(\Gamma_G(f)\bigr).
  \]

\item[\textup{\SGAInlineRefTarget{cor:sga3-VIB-7.3:item:iv}(iv)}]\phantomsection\label{concl:sga3-VIB-7.3-iv}
  \emph{If \(g\) is a \(k\)-point of \(G\), then the \(k\)-group
  \(\Gamma_G(\{g\})\) is commutative.}

\item[\textup{(v)}]\phantomsection\label{concl:sga3-VIB-7.3-v}
  \emph{Let \(A\) and \(B\) be two group subschemes of \(G\),
  geometrically reduced over \(k\). If \(A\) and \(B\) are
  invariant (respectively, characteristic), then so is
  \((A,B)\).}
\end{enumerate}

\emph{Proof.}
We first prove \textup{\SGARefLink{cor:sga3-VIB-7.3:item:i}{(i)}\SGARefLink{concl:sga3-VIB-7.3-i}{(a)}}. Put
\[
  \Gamma_S=\Gamma_G(f)_S.
\]
Then \(u^{-1}(\Gamma_S)\) is a closed \(S\)-group subscheme of
\(G_S\). By
\hyperref[concl:sga3-VIB-7.1-iii]{7.1\textup{(iii)}}, it is
therefore enough to show that \(f_S\) factors through
\(u^{-1}(\Gamma_S)\). Since
\[
  u\bigl(f_S(X_S)\bigr)\subset\Gamma_S
\]
and \(X_S\) is reduced, \(u\circ f_S\) factors through
\(\Gamma_S\); hence \(f_S\) factors through
\(u^{-1}(\Gamma_S)\). This proves \textup{\SGARefLink{concl:sga3-VIB-7.3-i-a}{(a)}}.

The first assertion of \textup{\SGARefLink{cor:sga3-VIB-7.3:item:b}{(b)}} follows by applying
\textup{\SGARefLink{concl:sga3-VIB-7.3-i-a}{(a)}} to \(u\) and \(u^{-1}\). In fact, it is enough to
assume
\[
  u\bigl(f_S(X_S)\bigr)\subset\Gamma_S
  \qquad\text{and}\qquad
  u^{-1}\bigl(f_S(X_S)\bigr)\subset\Gamma_S.
\]
The second assertion is the special case \(u=\operatorname{int}(g)\).

We now prove \textup{\SGARefLink{concl:sga3-VIB-7.3-i-c}{(c)}}. Consider the morphism of \(S\)-groups
\[
  G_S\longrightarrow G_S\times_SG_S,
  \qquad
  g\longmapsto\bigl(u(g),g\bigr),
\]
and let \(H\) be the inverse image of the diagonal. It is a group
subscheme of \(G_S\). Since \(G\) is separated over \(k\), by
\hyperref[sga3:VIA:0.3]{VI\(_A\), 0.3}, \(G_S\) is separated
over \(S\), and hence \(H\) is closed in \(G_S\). By hypothesis,
\(f_S\) factors through \(H\); therefore \(H\) contains
\(\Gamma_G(f)_S\), by
\hyperref[concl:sga3-VIB-7.1-iii]{7.1\textup{(iii)}}. This proves
the first assertion of \textup{\SGARefLink{cor:sga3-VIB-7.3:item:c}{(c)}}, and the second is the
special case \(u=\operatorname{int}(g)\).

We next prove \textup{\SGARefLink{cor:sga3-VIB-7.3:item:ii}{(ii)}}. Put
\[
  \Gamma=\Gamma_G(f)
\]
and let \(i:\Gamma\hookrightarrow G\) denote the inclusion. The
group \(H\) is contained in
\[
  N=\underline{\operatorname{Norm}}_G(\Gamma)
\]
if and only if, for every \(k\)-scheme \(S\) and every
\(h\in H(S)\),
\[
  \operatorname{int}(h)(\Gamma_S)\subset\Gamma_S.
\]
Applying this condition to \(h\) and \(h^{-1}\) shows that the
inclusion is then an equality. By \textup{\SGARefLink{cor:sga3-VIB-7.3:item:i}{(i)}\SGARefLink{concl:sga3-VIB-7.3-i}{(a)}}, this condition
holds if and only if
\[
  \operatorname{int}(h)\bigl(f_S(X_S)\bigr)\subset\Gamma_S.
\]
Similarly, \(H\) is contained in
\[
  C=\underline{\operatorname{Centr}}_G(\Gamma)
\]
if and only if, for every \(k\)-scheme \(S\) and every
\(h\in H(S)\),
\[
  \operatorname{int}(h)\circ i_S=i_S.
\]
By \textup{\SGARefLink{cor:sga3-VIB-7.3:item:i}{(i)}\SGARefLink{cor:sga3-VIB-7.3:item:c}{(c)}}, this holds if and only if
\[
  \operatorname{int}(h)\circ f_S=f_S.
\]
This proves \textup{\SGARefLink{cor:sga3-VIB-7.3:item:ii}{(ii)}}.

We now prove \textup{\SGARefLink{cor:sga3-VIB-7.3:item:iii}{(iii)}}. In view of \textup{\SGARefLink{cor:sga3-VIB-7.3:item:ii}{(ii)}}, the
hypothesis implies that \(\varphi\) factors through \(C\)
(respectively, \(N\)). The latter is a closed subgroup of \(G\),
by \hyperref[prop:sga3-VIB-6.2]{6.2}, and therefore contains
\(\Gamma_G(\varphi)\), by
\hyperref[concl:sga3-VIB-7.1-i]{7.1\textup{(i)}}.%
\begingroup
\renewcommand{\thefootnote}{\fnsymbol{footnote}}%
\footnote[1]{\label{note:sga3-VIB-7.3-Gamma-phi}Translator's note:
the French source prints \(\Gamma_G(f)\) here. Since
\(\varphi\), rather than \(f\), is the morphism shown to factor
through \(C\) (respectively, \(N\)), \SGARefLink{prop:sga3-VIB-7.1}{7.1}\textup{\SGARefLink{concl:sga3-VIB-7.1-i}{(i)}} gives
\(\Gamma_G(\varphi)\), as required by the statement.}
\endgroup

Assertion \textup{\SGARefLink{cor:sga3-VIB-7.3:item:iv}{(iv)}} follows from \textup{\SGARefLink{cor:sga3-VIB-7.3:item:iii}{(iii)}} by taking
both \(f\) and \(\varphi\) to be the closed immersion
\[
  \operatorname{Spec}k\hookrightarrow G
\]
defined by \(g\). Indeed, for every \(k\)-scheme \(S\), the sets
\(X(S)\) and \(Y(S)\) then each have a single element, which is
sent by \(f\), respectively by \(\varphi\), to \(g\).

Finally, we prove \textup{\SGARefLink{concl:sga3-VIB-7.3-iv}{(v)}}. Let
\[
  \nu:G\times_kG\longrightarrow G,
  \qquad
  (g,h)\longmapsto ghg^{-1}h^{-1},
\]
and let \(\nu'\) be its restriction to \(A\times_kB\). By
\hyperref[def:sga3-VIB-7.2.2]{7.2.2},
\[
  (A,B)=\Gamma_G(\nu').
\]
Let \(S\) be a \(k\)-scheme, and let \(u\) be an inner
automorphism, respectively an arbitrary \(S\)-group
automorphism, of \(G_S\). Then
\[
  u\circ\nu_S=\nu_S\circ(u\times_Su).
\]
The hypothesis implies that \(u\) induces automorphisms \(u_1\)
of \(A_S\) and \(u_2\) of \(B_S\). Consequently,
\phantomsection\label{form:sga3-VIB-7.3-v-commutator-image}%
\[
  u\bigl(\nu'_S(A_S\times_SB_S)\bigr)
    =\nu'_S\bigl(u_1(A_S)\times_Su_2(B_S)\bigr)
    =\nu'_S(A_S\times_SB_S).
\]
By \textup{\SGARefLink{concl:sga3-VIB-7.3-i}{(i)}\SGARefLink{concl:sga3-VIB-7.3-i-b}{(b)}}, \(u\) therefore induces an automorphism of
\((A,B)_S\). This proves \textup{\SGARefLink{concl:sga3-VIB-7.3-v}{(v)}}.
